% !TeX root = ../sections/02_exercises.tex

\begin{exercise}
	R-3. 環 $R$ 自身を左 $R$-加群とみたとき，加群としての同型
	\[
	\operatorname{End}_R(R)\cong R
	\]
	が成り立つことを示せ。
	さらに，これは環としての同型でもあることを示せ。
\end{exercise}

\begin{background}
	$R$ を左 $R$-加群とみるとは，$R$ の元 $a\in R$ が
	左からの積
	\[
	r\mapsto ar
	\]
	によって作用していると考えることである。
	
	$\operatorname{End}_R(R)$ は，左 $R$-加群としての
	$R$ から $R$ への準同型全体である。
	つまり，写像 $\varphi:R\to R$ が $\operatorname{End}_R(R)$ の元であるとは，
	任意の $r,s\in R$ に対して
	\[
	\varphi(r+s)=\varphi(r)+\varphi(s),
	\qquad
	\varphi(rs)=r\varphi(s)
	\]
	が成り立つことをいう。
	
	任意の $\varphi\in\operatorname{End}_R(R)$ は，$\varphi(1)$ によって完全に決まる。
	実際，
	\[
	\varphi(r)=\varphi(r\cdot 1)=r\varphi(1)
	\]
	である。
	
	ただし，$R$ が非可換環の場合には注意が必要である。
	左 $R$-加群 $R$ の自己準同型環は，一般には
	\[
	\operatorname{End}_R({}_RR)\cong R^{\operatorname{op}}
	\]
	となる。
	したがって，問題文のように環として
	\[
	\operatorname{End}_R(R)\cong R
	\]
	と述べるには，$R$ が可換環である，または積の向きについての約束を明確にする必要がある。
\end{background}

\begin{strategy}
	写像
	\[
	\Phi:\operatorname{End}_R(R)\longrightarrow R
	\]
	を
	\[
	\Phi(\varphi)=\varphi(1)
	\]
	で定める。
	
	まず，$\Phi$ が加法について準同型であることを確認する。
	次に，任意の $a\in R$ に対して
	\[
	\varphi_a(r)=ra
	\]
	と定めると，$\varphi_a$ は左 $R$-加群準同型であり，
	\[
	\Phi(\varphi_a)=a
	\]
	となることを確認する。
	これにより $\Phi$ は全射である。
	
	さらに，
	\[
	\Phi(\varphi)=0
	\]
	ならば $\varphi(1)=0$ であり，任意の $r\in R$ について
	\[
	\varphi(r)=r\varphi(1)=0
	\]
	となるので，$\varphi=0$ である。
	したがって $\Phi$ は単射である。
	
	環構造については，合成の順序に注意する。
	$\varphi_a(r)=ra$，$\varphi_b(r)=rb$ とすると，
	\[
	(\varphi_a\circ\varphi_b)(r)
	=
	\varphi_a(rb)
	=
	(rb)a
	=
	r(ba)
	\]
	である。
	したがって，
	\[
	\varphi_a\circ\varphi_b=\varphi_{ba}
	\]
	となる。
	
	このため，一般の非可換環では
	\[
	\operatorname{End}_R({}_RR)\cong R^{\operatorname{op}}
	\]
	である。
	特に $R$ が可換環ならば $R^{\operatorname{op}}=R$ なので，
	環としても
	\[
	\operatorname{End}_R(R)\cong R
	\]
	が成り立つ。
\end{strategy}